% CLAUDE REPORT::BRIDGE-CALIBRATION-0001::2026-11-10T23:51-08:00
% !TEX program = xelatex
% Compiles with XeLaTeX or pdfLaTeX (ASCII source).
\documentclass[11pt]{article}
\usepackage[margin=1in]{geometry}
\usepackage{amsmath,amssymb,amsthm}
\usepackage{xurl}
\newtheorem{theorem}{Theorem}
\newtheorem{proposition}[theorem]{Proposition}
\newtheorem{corollary}[theorem]{Corollary}
\newtheorem{lemma}[theorem]{Lemma}
\theoremstyle{definition}
\newtheorem{definition}[theorem]{Definition}
\theoremstyle{remark}
\newtheorem*{remark}{Remark}
\newcommand{\SAT}{\mathrm{SAT}}
\newcommand{\UNSAT}{\mathrm{UNSAT}}
\newcommand{\Wbar}{\overline W}
\newcommand{\poly}{\mathrm{poly}}
\newcommand{\status}[1]{\textbf{[#1]}}

\title{CLAUDE BRIDGE-CALIBRATION-0001\\[2pt]
\large How strong is the ``weakest sufficient bridge''?\\
Exact collapse of $\SAT\in P\Rightarrow \Wbar_R(N)\le\poly(N)$, its monotonicity in $R$,\\
and the feasible-provability form for Extended Frege}
\author{CLAUDE scheduled routine (independent of the five LSPV roles)}
\date{2026-11-10 23:51 America/Los\_Angeles}

\begin{document}
\maketitle

\noindent\textbf{Status.} \textbf{P versus NP remains OPEN.} This is an additive
report of the scheduled CLAUDE routine working on the shared worst-case
minimum-derivation directive. It changes no LSPV role, schedule, allocation,
ownership region, historical rule grammar or weight, and it does not edit root
\texttt{main.tex}, \texttt{proof1.txt} or \texttt{proof2.txt}. Living log:
\path{research/agents/claude/CLAUDE_LIVING_LOG.md}.

\section{Shared objective, restated}

The user's multisign DAG hypergraph view of a 3-CNF $F$ is its signed clause
hypergraph; a \emph{reduction sequence} is a derivation in a fixed sound rule
system $R$; for $F\in\UNSAT$ the minimal reduced logical state is the empty
edge $\bot$. With unit or bounded fixed weights, the minimum work is
\[
W_R(F)=\min_{\pi:F\vdash_R\bot}\operatorname{cost}_R(\pi),\qquad
\Wbar_R(N)=\max_{F\in\UNSAT,\ |\langle F\rangle|\le N} W_R(F).
\]
Granting free knowledge of the optimal sequence is exactly the nondeterministic
\emph{proof-size} measure of $R$, so the user's remark is correct: the
optimum is a legitimate lower bound for any procedure whose execution trace
\emph{is} an $R$-derivation (for example DPLL/CDCL traces are Resolution
refutations), and search can only add work. Write
\[
\mathrm{LB}_R:\ \Wbar_R\ \text{is superpolynomial},\qquad
\mathrm{B3}_R:\ \SAT\in P\Longrightarrow \Wbar_R(N)\le\poly(N).
\]
(The project trackers call these \texttt{[P2-B1]} and \texttt{[P2-B2]} for
$R=R_{\mathrm{mix341}}$; \texttt{[P2-B4]} records $\mathrm{LB}_R\wedge\mathrm{B3}_R\Rightarrow P\ne NP$.)

The directive asks agents to \emph{prioritize the weakest sufficient bridge}
$\mathrm{B3}_R$. This report calibrates exactly how weak that bridge can be.

\section{Main results}

Throughout, $R$ is a sound and complete refutation system for CNF whose
cost measure is polynomially related to proof bit-length and whose proofs are
checkable in polynomial time (a Cook--Reckhow system). All statements below
are about one fixed such $R$.

\begin{theorem}[Bridge collapse; \status{PROVED/EXACT}]\label{thm:collapse}
For every fixed $R$,
\[
\mathrm{B3}_R\iff\bigl(\SAT\notin P\ \vee\ \Wbar_R\le\poly\bigr),
\qquad
\bigl(\mathrm{LB}_R\wedge\mathrm{B3}_R\bigr)\iff\bigl(\mathrm{LB}_R\wedge P\neq NP\bigr).
\]
\end{theorem}

\begin{proof}
The first equivalence is the definition of material implication. For the
second: ($\Rightarrow$) is \texttt{[P2-B4]}: if $\SAT\in P$ then
$\mathrm{B3}_R$ gives $\Wbar_R\le\poly$, contradicting $\mathrm{LB}_R$.
($\Leftarrow$) If $P\ne NP$ then $\SAT\notin P$, so $\mathrm{B3}_R$ holds
vacuously by the first equivalence.
\end{proof}

\begin{corollary}[Known-lower-bound systems; \status{PROVED/EXACT from KNOWN THEOREMS}]\label{cor:known}
If $\mathrm{LB}_R$ is a theorem, then $\mathrm{B3}_R\iff P\ne NP$ (provably in
ZFC). This applies to Resolution (Haken 1985), Cutting Planes (Pudl\'ak 1997),
bounded-depth Frege (Ajtai 1988; Kraj\'i\v{c}ek--Pudl\'ak--Woods 1995;
Pitassi--Beame--Impagliazzo 1993), Polynomial Calculus
(Razborov 1998; Ben-Sasson--Impagliazzo 1999) and the artificial
$R_{\mathrm{clock}}$ of \texttt{[P2-B6]}.
\end{corollary}

So for every system in which the first priority (a superpolynomial worst-case
minimum) is already established, the ``weakest sufficient bridge'' \emph{is}
the full P versus NP problem; there is no slack between them. This unifies two
earlier project results as special cases: \texttt{[P2-B6]} (LSPV-1/462,
$R_{\mathrm{clock}}$) and the TQBF-oracle Haken firewall in root
\texttt{main.tex} (oracle-relative B1/B2/B3 fail for Resolution).

\begin{corollary}[Generic-bridge test; \status{PROVED/EXACT}]\label{cor:generic}
Any argument for $\mathrm{B3}_R$ that uses only soundness, completeness and
polynomial-time checkability of $R$ also proves $\mathrm{B3}_{R_{\mathrm{clock}}}$, hence
proves $P\ne NP$ outright. A noncircular bridge must therefore use rule-specific
\emph{simulating power} of $R$ (the direction formalized in LSPV-5/409:
$\mathrm{B3}_R\iff Q_A\le_W R$ for the identity system $Q_A$ of a hypothetical
polynomial decider $A$).
\end{corollary}

\begin{proposition}[Monotonicity in proof strength; \status{PROVED/EXACT}]\label{prop:mono}
Suppose $R'$ $W$-simulates $R$: there is a polynomial $p$ with
$W_{R'}(F)\le p(|\langle F\rangle|+W_R(F))$ for all $F\in\UNSAT$. Then
$\mathrm{B3}_R\Rightarrow\mathrm{B3}_{R'}$ and $\mathrm{LB}_{R'}\Rightarrow\mathrm{LB}_R$.
\end{proposition}

\begin{proof}
$\Wbar_{R'}(N)\le p(N+\Wbar_R(N))$, so a polynomial bound for $R$ transfers to $R'$;
take contrapositives for the lower bound.
\end{proof}

Thus the two obligations move in opposite directions as $R$ strengthens:
the bridge gets easier, the lower bound gets harder, and by
Theorem~\ref{thm:collapse} their conjunction is always exactly
$\mathrm{LB}_R\wedge P\ne NP$. At the top end:

\begin{proposition}[p-optimal end; \status{PROVED/EXACT from KNOWN THEOREMS}]\label{prop:opt}
If $R$ $W$-simulates every Cook--Reckhow system, then $\mathrm{B3}_R$ holds
unconditionally and $\mathrm{LB}_R\iff NP\ne coNP$.
\end{proposition}

\begin{proof}
If $\SAT\in P$ via $A$, the identity system $Q_A$ (LSPV-5/409) is polynomially
bounded and $R$ simulates it. For the second claim use Cook--Reckhow (1979):
$NP=coNP$ iff some system is polynomially bounded, and $R$ simulates it.
\end{proof}

Whether such an $R$ exists is the Kraj\'i\v{c}ek--Pudl\'ak optimal proof
system problem (\status{OPEN}; $E=NE$ suffices, Kraj\'i\v{c}ek--Pudl\'ak 1989;
Egidy--Gla\ss{}er, STOC 2025, give oracles showing positive answers cannot
relativize). This is the size-optimality route already recorded as LSPV-1/336
and LSPV-5/296; it is credited, not claimed new.

\subsection{The intermediate point: Extended Frege and feasible provability}

\begin{theorem}[Feasible-provability bridge for EF; \status{KNOWN THEOREM}]\label{thm:ef}
Let $A$ be a polynomial-time algorithm, given as a $\mathrm{PV}$ function symbol,
and suppose $\mathrm{PV}\vdash \forall F\,\forall a\,
\bigl(A(F)=\UNSAT\rightarrow \neg\mathrm{Sat}(F,a)\bigr)$. Then for every $F$
with $A(F)=\UNSAT$ there is an EF refutation of $F$ of size
$\poly(|\langle F\rangle|)$, constructible in polynomial time.
Consequently, if $\SAT\in P$ via such a PV-provably sound $A$, then
$\Wbar_{EF}\le\poly$.
\end{theorem}

\begin{proof}[Proof sketch]
Cook (1975): PV-provable universal statements have polynomial-size,
polynomial-time constructible EF proofs of their propositional translations.
Instantiate the translation at length $|\langle F\rangle|$, substitute the
constant bits of $F$, prove the variable-free evaluation $A(F)=\UNSAT$ in EF with
one extension variable per circuit gate, apply modus ponens, and use the
polynomial-size EF equivalence between $\mathrm{Sat}(F,a)$ at constant $F$ and
$F(a)$. The result is an EF derivation of $\neg F$.
\end{proof}

Cook and Kraj\'i\v{c}ek (JSL 2007) sharpen this:
$\mathrm{PV}\vdash NP=coNP$ iff PV defines a polynomial-time algorithm that
outputs either a satisfying assignment or an EF refutation; the same holds for
$S^1_2$. Byd\v{z}ovsk\'y--Kraj\'i\v{c}ek--Oliveira (arXiv 1905.12935) note that
``no polynomial-time $f$ is PV-provably a SAT solver'' follows from EF not being
polynomially bounded.

\begin{corollary}[Exact residual gap for $R=EF$ and its extensions; \status{PROVED/EXACT relative to Theorem~\ref{thm:ef}}]
For every $R$ that $W$-simulates EF,
\[
\mathrm{H}_{\mathrm{feas}}:\ \SAT\in P\Rightarrow\exists A\ \text{polytime with PV-provable soundness of its UNSAT answers}
\]
implies $\mathrm{B3}_R$. Hence $\mathrm{LB}_R\wedge\mathrm{H}_{\mathrm{feas}}\Rightarrow P\ne NP$.
Unconditionally, $\mathrm{LB}_{R}$ alone already implies that no polynomial-time
SAT algorithm is PV-provably sound.
\end{corollary}

$\mathrm{H}_{\mathrm{feas}}$ is \status{OPEN}; it is a statement about the
\emph{provability} of a hypothetical algorithm's correctness, and nothing known
derives it from $\SAT\in P$. It is nonetheless the sharpest known
rule-specific bridge, and it is \emph{not} equivalent to $P\ne NP$ in the sense
of Corollary~\ref{cor:known}, because $\mathrm{LB}_{EF}$ is not known.

\section{Consequence for $R_{\mathrm{mix341}}$ (\status{CONDITIONAL})}

The corpus records that $R_{\mathrm{mix341}}$ contains paid auxiliary
definitions and has polynomial mixed refutations of pigeonhole formulas
(LSPV-3/346; Cook ER countercontrol in LSPV-1/371). \emph{If} its fully
authenticated grammar \texttt{[P2-D1]} $W$-simulates Extended Resolution
(equivalently EF), then:
\begin{enumerate}
\item \texttt{[P2-B1]} would imply a superpolynomial EF lower bound, a central
open problem of proof complexity for which no superpolynomial lower bound is
known for any explicit family, and which barrier results (Razborov's
pseudorandom-generator hardness conjectures; Kraj\'i\v{c}ek's proof-complexity
generators) suggest needs new methods;
\item \texttt{[P2-B2]} would follow from $\mathrm{H}_{\mathrm{feas}}$
by Theorem~\ref{thm:ef} and Proposition~\ref{prop:mono}.
\end{enumerate}
If instead $R_{\mathrm{mix341}}$ is strictly weaker than ER and a superpolynomial
lower bound for it is ever proved, Corollary~\ref{cor:known} makes
\texttt{[P2-B2]} equivalent to $P\ne NP$. Either way, deciding whether the
authenticated $R_{\mathrm{mix341}}$ simulates ER is the single most useful
classification question for the Route B spine: it tells the project which of
the two obligations is the frontier. \status{OPEN}.

\section{Hard-instance survival ledger (task item 4)}

\begin{center}
\small
\begin{tabular}{llll}
Family & Hard for (KNOWN) & Collapses in (KNOWN) & Status in strong $R$\\\hline
Pigeonhole PHP & Res (Haken), bd-depth Frege & CP, ER/EF (Cook 1976; Buss) & dead for any $R\ge$ER\\
Tseitin on expanders & Res (Urquhart; Ben-Sasson--Wigderson) & GF(2) / Res(lin), PC over $\mathbb F_2$ & dead for $R$ with XOR\\
random 3-CNF & Res (Chv\'atal--Szemer\'edi), PC, CP & none known & candidate; EF \status{OPEN}\\
Proof-complexity generators & conjectured for EF (Kraj\'i\v{c}ek) & none known & \status{CANDIDATE} only
\end{tabular}
\end{center}

\noindent Families that are hard only for systems whose lower bound is already
known contribute, by Corollary~\ref{cor:known}, nothing toward the bridge.
A family is useful to Route B only if it is a candidate against a system at or
above the bridge's simulating strength, which today means EF-hardness
candidates.

\section{Barrier-transfer literature preflight}

\textbf{Obstacles.} (O1) $\mathrm{B3}_R$ for a fixed $R$; (O2)
$\mathrm{LB}_R$ for any $R$ strong enough to make O1 rule-specific.

\textbf{Literature checked this run.} Cook--Kraj\'i\v{c}ek 2007 (PV and
$NP=coNP$); Byd\v{z}ovsk\'y--Kraj\'i\v{c}ek--Oliveira, arXiv 1905.12935;
Egidy--Gla\ss{}er STOC 2025 and Egidy arXiv 2602.02294 (oracles against
relativizable positive answers to optimal-proof-system questions);
Lu--Santhanam--Tzameret, ECCC TR25-134 / ITCS 2026 (AC$^0[p]$-Frege lower
bounds for DNFs; constant-depth algebraic frontier);
\emph{The Proof Analysis Problem}, arXiv 2506.16956 (FOCS 2025: EF proof analysis
NP-complete; EF formalizes NP-hardness of automating Resolution);
PC with extension rule, Theory of Computing vol.~22 (2026) (algebraic systems that
simulate EF). None proves an EF lower bound or $\mathrm{H}_{\mathrm{feas}}$.

\textbf{Transfer classification.}
(i) Bounded arithmetic $\to$ proof systems (Cook translation): \emph{directly
transferable theorem} (Theorem~\ref{thm:ef}).
(ii) Control-theory barrier certificates $\leftrightarrow$ refutations: a
refutation is a certificate that no state reaches ``all clauses satisfied'';
\emph{structural analogy} only, since certificate size bounds there are not
worst-case lower bounds.
(iii) VLSI area--time and communication-cut lower bounds $\to$ proof size:
\emph{directly transferable} only for weak systems (CP via monotone
interpolation and communication, Pudl\'ak 1997); for EF it is a \emph{failed
transfer} under standard cryptographic assumptions (Kraj\'i\v{c}ek--Pudl\'ak 1998:
EF has no feasible interpolation if RSA is secure).

\textbf{Post-quantum ledger.} Item 5 (worst-case/average-case architecture) and
item 6 (special soundness) are \emph{failed transfers} for O1 here: they would
extract proofs from a black-box solver, which Corollary~\ref{cor:generic}
shows must use rule-specific simulation of $Q_A$ that a black box does not
supply. Items 1--4 do not bear on O1/O2 in this report.
Relativization: Theorem~\ref{thm:collapse} relativizes trivially; it is a
logical calibration, not a separation technique.

\section{Corpus reconciliation}

\paragraph{CORPUS REVIEWED.} Repository tree; \texttt{AGENTS.md}; root
\texttt{main.tex} (headers, LSPV5 REFEREE294--297 executive blocks, LSPV-1
336/347/350/351/371/373/374/376 entries, TQBF--Haken firewall);
\texttt{proof1.txt} header and target nodes; \texttt{proof2.txt} nodes
\texttt{[P2-D1]}--\texttt{[P2-D3]}, \texttt{[P2-B1]}--\texttt{[P2-B9]},
LSPV-5/409, LSPV-1/466, LSPV-1/462 (\texttt{[P2-B6]}); git log through
LSPV-3/563; earlier CLAUDE reports (world ledger, TPORT, ONETAPE).
The five-role knot/link handoff (2026-11-09) was noted but not acted on: it is
addressed to the LSPV roles.

\paragraph{DUPLICATION / CONTRADICTION AUDIT.}
Already in corpus and credited, not claimed: B1/B4 logic (\texttt{[P2-B4]});
size-optimality $\Rightarrow$ B3 (LSPV-1/336, LSPV-5/296); relativization barrier
for unchanged $R$ (LSPV-5/296); $\mathrm{B3}_R\iff Q_A\le_W R$ (LSPV-5/409);
$R_{\mathrm{clock}}$ (\texttt{[P2-B6]}); oracle Haken firewall.
New here: the general collapse Theorem~\ref{thm:collapse} and
Corollary~\ref{cor:known} naming every system with a known lower bound, the
generic-bridge test (Corollary~\ref{cor:generic}), the monotonicity
Proposition~\ref{prop:mono}, and the explicit import of the Cook /
Cook--Kraj\'i\v{c}ek feasible-provability bridge with gap $\mathrm{H}_{\mathrm{feas}}$.
A repository search found no earlier PV/Cook-1975/feasible-provability entry in
\texttt{proof2.txt}. No contradiction with any recorded status was found.

\paragraph{WHAT CHANGED RELATIVE TO ROOT MAIN.} Nothing in root
\texttt{main.tex} was edited (this routine owns no region). Suggested for the
LSPV-5 arbiter: add Theorem~\ref{thm:collapse} as a status note under
\texttt{[P2-B2]}, and add ``does authenticated $R_{\mathrm{mix341}}$ $W$-simulate
ER?'' as a classification target under \texttt{[P2-D1]}.

\paragraph{ADOPTED / REJECTED / RETIRED PRIOR COMPONENTS.}
Adopted: \texttt{[P2-D2]} definitions; LSPV-5/409 simulation equivalence.
Rejected: none. Retired: none. The directive's description of B3 as a
``weakest'' bridge is retained, with the precision that for every system with a
known lower bound it is \emph{equivalent} to the target.

\paragraph{STALE OR OVERSTATED CLAIMS FOUND.} None in the trackers. Note for
future wording: ``prove the weakest sufficient bridge first'' must not be read
as an easier subgoal when $R$ is Resolution, CP, PC or bounded-depth Frege.

\section{Recommendations for the five roles}

\begin{enumerate}
\item \textbf{General SAT Lower-Bound Lead:} adopt Theorem~\ref{thm:collapse} as
the formal frame; target either $\mathrm{H}_{\mathrm{feas}}$-style rule-specific
bridges for systems at or above EF, or decide the ER-simulation question for
$R_{\mathrm{mix341}}$.
\item \textbf{SAT-Circuit Reduction Builder:} test whether the authenticated DEF
rule gives a $W$-simulation of Extended Resolution (polynomial overhead per
extension axiom).
\item \textbf{Hypergraph \& Proof Certificates:} prefer EF-hardness candidates
(random 3-CNF, proof-complexity generator formulas) over PHP and Tseitin, which
are dead in any $R\ge$ ER with XOR.
\item \textbf{Time--Space Resource Auditor:} confirm that $R_{\mathrm{mix341}}$ cost
is polynomially related to proof bit-length, the hypothesis needed for
Proposition~\ref{prop:opt} and Theorem~\ref{thm:ef}.
\item \textbf{Independent Proof Red Team:} apply Corollary~\ref{cor:generic}
to every proposed B3 argument: if it would also apply to
$R_{\mathrm{clock}}$, it is either a full P$\ne$NP proof or wrong.
\end{enumerate}

\section*{Status summary}
\begin{tabular}{ll}
Theorem~\ref{thm:collapse}, Corollaries~\ref{cor:known},~\ref{cor:generic}, Prop.~\ref{prop:mono} & PROVED/EXACT\\
Proposition~\ref{prop:opt} & PROVED/EXACT (existence of such $R$ OPEN)\\
Theorem~\ref{thm:ef} & KNOWN THEOREM\\
$\mathrm{H}_{\mathrm{feas}}$; ER-simulation by $R_{\mathrm{mix341}}$ & OPEN\\
\texttt{[P2-B1]}, \texttt{[P2-B2]}, P versus NP & OPEN
\end{tabular}

\medskip\noindent One OPEN load-bearing arrow means $P\ne NP$ is not proved.
\textbf{P versus NP remains OPEN.}

\end{document}
% END_OF_VERIFIED_REPORT::CLAUDE::BRIDGE-CALIBRATION-0001
